
\documentclass[11pt]{article}

\usepackage{amsmath,amssymb,amsthm}
\usepackage{geometry}
\usepackage{hyperref}
\usepackage{physics}
\usepackage{bm}
\usepackage{mathrsfs}
\usepackage{enumitem}

\geometry{margin=1in}

\title{\textbf{Magnetism, Relativity, and the Geometry of Phase Space:\
Why Classical Electrodynamics Cannot Support Magnetic Matter}}

\author{}
\date{}

\begin{document}

\maketitle

\begin{abstract}
It is a classical but conceptually profound result that classical electrodynamics, even when combined with classical statistical mechanics, predicts the complete absence of both diamagnetism and paramagnetism in equilibrium. This review presents a unified geometric and relativistic analysis of this fact. We show that the obstruction to magnetism is not dynamical but structural, rooted in the symplectic geometry of classical phase space, the nature of gauge coupling, and the constraints imposed by relativistic covariance. We connect the Bohr–van Leeuwen theorem to the Ehlers–Pirani–Schild (EPS) construction of spacetime geometry and to the no–interaction theorems of relativistic dynamics. The analysis demonstrates that magnetism is not a weakly quantum phenomenon but an intrinsically quantum one, inseparable from noncommutativity, spin, and spectral discreteness.
\end{abstract}

\tableofcontents

%------------------------------------------------

\section{Introduction}

Magnetism is among the most familiar manifestations of matter–field interaction, yet its theoretical status is remarkably subtle. Classical electrodynamics describes magnetic forces with perfect accuracy at the level of equations of motion, and yet—when embedded into classical statistical mechanics—it predicts that matter in thermal equilibrium exhibits \emph{no magnetic response whatsoever}.

This apparent paradox was resolved in the early twentieth century by the Bohr–van Leeuwen theorem, which shows that magnetization vanishes identically in classical equilibrium systems. Far from being a technical curiosity, this result reveals a deep incompatibility between classical mechanics and magnetism as a thermodynamic phenomenon.

In this review we argue that this incompatibility is structural and geometric. Magnetism is excluded by the combined requirements of:
\begin{itemize}
\item canonical phase space structure,
\item gauge invariance,
\item relativistic covariance,
\item and the absence of intrinsic internal degrees of freedom.
\end{itemize}

We further show that the obstruction is closely related to the no–interaction theorem and to the geometric foundations of spacetime articulated by Ehlers, Pirani, and Schild.

---

\section{Classical Electrodynamics as a Gauge Theory}

\subsection{Minimal Coupling and Canonical Structure}

For a classical charged particle,
\begin{equation}
H = \frac{1}{2m}\left( \mathbf{p} - \frac{q}{c}\mathbf{A}(\mathbf{x}) \right)^2 + q\phi(\mathbf{x}),
\end{equation}
with symplectic form
\begin{equation}
\omega = d\mathbf{p} \wedge d\mathbf{x}.
\end{equation}

The electromagnetic field enters only through a shift of canonical momentum. This is crucial: the map
\[
\mathbf{p} \mapsto \mathbf{p} - \frac{q}{c}\mathbf{A}
\]
is a canonical transformation.

Hence electromagnetism does not deform phase space—it merely relabels it.

---

\subsection{Gauge Fields as Connections}

The vector potential $\mathbf{A}$ defines a connection on a principal $U(1)$ bundle over configuration space. In classical mechanics, however, this connection has no physical holonomy: only the field strength enters the equations of motion.

This stands in sharp contrast to quantum mechanics, where the connection acts on wavefunctions and produces observable holonomy (Aharonov–Bohm effect).

---

\section{The Bohr–van Leeuwen Theorem Revisited}

The canonical partition function,
\[
Z = \int d^{3N}x, d^{3N}p ; e^{-\beta H},
\]
is invariant under the canonical momentum shift induced by the vector potential.

Thus,
\[
\frac{\partial Z}{\partial \mathbf{B}} = 0,
\qquad
\mathbf{M} = 0.
\]

This result is entirely independent of interaction details, boundary conditions, or particle statistics.

The theorem reflects not a cancellation of forces but a deeper invariance of the Liouville measure.

---

\section{Geometric Interpretation: Phase Space and Symplectic Flatness}

The phase space $(\Gamma,\omega)$ of classical mechanics is symplectically flat. Electromagnetic coupling does not introduce curvature into $\omega$, only a gauge-dependent reparametrization.

Thus:
\begin{itemize}
\item No topological obstruction exists.
\item No geometric phase is generated.
\item No preferred orientation arises in equilibrium.
\end{itemize}

In modern language, magnetism requires curvature in the \emph{quantum} line bundle, not merely in spacetime.

---

\section{Relativistic Covariance and the Ehlers--Pirani--Schild Framework}

Ehlers, Pirani, and Schild (EPS) showed that spacetime geometry can be reconstructed from the behavior of light rays (conformal structure) and free-fall trajectories (projective structure).

In this framework:
\begin{itemize}
\item The conformal structure encodes electromagnetism only indirectly.
\item The projective structure encodes inertial motion.
\end{itemize}

Classical charged particles do not modify these structures; they merely respond to them.

Thus electromagnetic interactions do not alter the spacetime geometry in the same sense as gravitation.

This sharply limits the role electromagnetism can play in defining invariant physical structure, reinforcing the absence of intrinsic magnetization.

---

\section{Connection to the No–Interaction Theorem}

The Currie–Jordan–Sudarshan theorem states that relativistic systems of particles with canonical phase space and Poincaré covariance admit only trivial interactions.

Magnetism would require velocity-dependent, nonlocal interactions violating these assumptions.

Thus:
\begin{itemize}
\item Classical electromagnetism is compatible with relativity only as an external field theory.
\item True interacting many-body dynamics require abandoning canonical particle mechanics.
\end{itemize}

This aligns perfectly with the Bohr–van Leeuwen result: magnetism cannot arise within the allowed classical structures.

---

\section{Why Semiclassical Fixes Fail}

Attempts to restore magnetism via semiclassical quantization fail because:

\begin{itemize}
\item Action–angle quantization preserves classical phase space structure.
\item Gauge potentials remain removable.
\item Spin is absent.
\end{itemize}

Only full quantization introduces noncommuting observables and topologically nontrivial bundles.

---

\section{Quantum Resolution}

Quantum mechanics alters the ontology:

\begin{itemize}
\item $[x_i,p_j]=i\hbar\delta_{ij}$ destroys canonical equivalence.
\item Magnetic fields discretize spectra (Landau levels).
\item Spin introduces intrinsic magnetic moments.
\item Berry curvature endows phase space with topology.
\end{itemize}

Magnetism emerges as a geometric quantum phenomenon.

---

\section{Conceptual Synthesis}

Magnetism fails classically because:
\begin{enumerate}
\item Classical phase space is flat.
\item Gauge fields are pure connections.
\item Relativistic covariance forbids interacting particle dynamics.
\item Statistical mechanics averages over canonical symmetries.
\end{enumerate}

Quantum theory breaks all four assumptions.

---

\section{Conclusion}

The absence of diamagnetism and paramagnetism in classical electrodynamics is not a technical anomaly but a structural necessity. When viewed through the lenses of symplectic geometry, relativistic covariance, and gauge theory, magnetism appears as an intrinsically quantum phenomenon, inseparable from noncommutativity and topological structure.

In this sense, magnetism plays a role analogous to spacetime curvature in gravity: it signals the breakdown of classical kinematics and the necessity of a deeper geometric framework.

---

\section*{References}

\begin{enumerate}
\item N. Bohr, \emph{Studier over Metallernes Elektrontheori}, Copenhagen (1911).
\item J. H. van Leeuwen, J. Phys. Radium \textbf{2}, 361 (1921).
\item J. Ehlers, F. A. E. Pirani, A. Schild, \emph{The Geometry of Free Fall and Light Propagation}, in \emph{General Relativity}, Clarendon (1972).
\item H. Goldstein, C. Poole, J. Safko, \emph{Classical Mechanics}.
\item L. D. Landau, E. M. Lifshitz, \emph{Statistical Physics}.
\item R. Haag, \emph{Local Quantum Physics}.
\item J. C. Taylor, \emph{Gauge Theories of Weak Interactions}.
\item M. Stone, \emph{The Physics of Quantum Fields}.
\end{enumerate}

\end{document}
